\Needspace{8\baselineskip}
\Article{256}
Die folgenden allgemeinen Bemerkungen werden diese Lösung noch deutlicher machen.

I. Ist die Ordnung \(O\) von einer eigentlich primitiven Ordnung abgeleitet, so teilt \(AA\) die Zahl \(D\); ist dagegen \(O\) uneigentlich primitiv oder von einer uneigentlich primitiven Ordnung abgeleitet, so ist \(A\) gerade, \(D\) durch \(\frac14 AA\) teilbar und der Quotient \(\equiv1\ (\mathrm{mod.}\,4)\). Somit teilt das Quadrat jedes Teilers von \(A\) entweder \(D\) selbst oder wenigstens \(4D\), und im letzteren Fall ist der Quotient stets \(\equiv1\ (\mathrm{mod.}\,4)\).

II. Teilt \(aa\) die Zahl \(D\), so sind alle Werte des Ausdrucks \(\surd D\ (\mathrm{mod.}\,aa)\), die zwischen \(0\) und \(aa-1\) liegen, gleich \(0,a,2a\ldots aa-a\), und \(a\) ist daher die Anzahl der Formen in \(V\); eigentlich primitiv sind darunter aber nur so viele, wie es unter den Zahlen
\[
\frac{D}{aa},\quad\frac{D}{aa}-1,\quad\frac{D}{aa}-4\ldots\frac{D}{aa}-(a-1)^2
\]
\SourceLeaf{285}
solche gibt, die mit \(a\) keinen gemeinsamen Teiler haben. Ist \(a=1\), so besteht \(V\) aus einer einzigen Form, \((1,0,-D)\), die stets eigentlich primitiv ist. Ist \(a\) gleich \(2\) oder irgendeiner Potenz von \(2\), so ist die Hälfte jener \(a\) Zahlen gerade, die andere ungerade; daher gibt es in \(V\) \(\frac12 a\) eigentlich primitive Formen. Ist \(a\) eine andere Primzahl \(p\) oder eine Potenz der Primzahl \(p\), so sind drei Fälle zu unterscheiden: Alle jene \(a\) Zahlen sind zu \(a\) teilerfremd und somit alle Formen in \(V\) eigentlich primitiv, wenn \(\frac{D}{aa}\) nicht durch \(p\) teilbar und zugleich quadratischer Nichtrest modulo \(p\) ist; teilt dagegen \(p\) die Zahl \(\frac{D}{aa}\), so gibt es in \(V\) \(\frac{(p-1)a}{p}\) eigentlich primitive Formen; ist schließlich \(\frac{D}{aa}\) quadratischer Rest modulo \(p\) und nicht durch \(p\) teilbar, so gibt es in \(V\) \(\frac{(p-2)a}{p}\) eigentlich primitive Formen. Das alles lässt sich ohne Mühe beweisen. Setzt man allgemein \(a=2^\nu p^\pi q^\chi r^\rho\ldots\), wobei \(p,q,r\) usw. verschiedene ungerade Primzahlen bezeichnen, so ist die Anzahl der eigentlich primitiven Formen in \(V\) gleich \(NPQR\ldots\), wobei man setzen muss
\[
\begin{aligned}
N&=1\quad(\text{wenn }\nu=0)\text{ oder }N=2^{\nu-1}\ (\text{wenn }\nu>0)\\
P&=p^\pi\quad(\text{wenn }\tfrac{D}{aa}\text{ quadratischer Nichtrest modulo }p\text{ ist})\text{ oder}\\
P&=(p-1)p^{\pi-1}\quad(\text{wenn }\tfrac{D}{aa}\text{ durch }p\text{ teilbar ist})\text{ oder}\\
P&=(p-2)p^{\pi-1}\quad(\text{wenn }\tfrac{D}{aa}\text{ quadratischer Rest modulo }p\\ &\qquad\text{und nicht durch }p\text{ teilbar ist})
\end{aligned}
\]
und \(Q,R\) usw. aus \(q,r\) usw. ebenso bestimmen muss wie \(P\) aus \(p\).

III. Teilt \(aa\) die Zahl \(D\) nicht, so ist \(\frac{4D}{aa}\) ganzzahlig und \(\equiv1\ (\mathrm{mod.}\,4)\), und die Werte des Ausdrucks \(\surd D\ (\mathrm{mod.}\,aa)\) sind \(\frac12 a,\frac32 a,\frac52 a\ldots aa-\frac12 a\). Somit ist die Anzahl der Formen in \(V\) gleich \(a\); eigentlich primitiv sind darunter aber nur so viele, wie es unter den Zahlen
\[
\frac{D}{aa}-\tfrac14,\quad\frac{D}{aa}-\tfrac94,\quad\frac{D}{aa}-\tfrac{25}{4}\quad\ldots\quad\frac{D}{aa}-(a-\tfrac12)^2
\]
solche gibt, die zu \(a\) teilerfremd sind. Gilt \(\frac{4D}{aa}\equiv1\ (\mathrm{mod.}\,8)\), so sind alle diese Zahlen gerade, und es gibt daher in \(V\) keine eigentlich primitive Form; gilt dagegen \(\frac{4D}{aa}\equiv5\ (\mathrm{mod.}\,8)\). so sind alle jene Zahlen ungerade und somit alle Formen in \(V\) eigentlich primitiv, wenn \(a\) gleich \(2\) oder einer Potenz von \(2\) ist. Allgemein gibt es in diesem Fall in \(V\) so viele eigentlich primitive Formen, wie unter jenen Zahlen durch keinen ungeraden Primteiler von \(a\) teilbar sind. Diese Anzahl ist \(NPQR\ldots\), wenn \(a=2^\nu p^\pi q^\chi r^\rho\ldots\), wobei man \(N=2^\nu\) setzen und \(P,Q,R\) usw. aus \(p,q,r\) usw. ebenso ableiten muss wie im vorigen Fall.

\SourceLeaf{286}
IV. Auf diese Weise lassen sich also die Anzahlen der eigentlich primitiven Formen in \(V,V',V''\) usw. bestimmen. Für die Summe all dieser Anzahlen findet man ohne große Mühe die folgende allgemeine Regel: Ist \(A=2^\nu\mathfrak{A}^\alpha\mathfrak{B}^\beta\mathfrak{C}^\gamma\ldots\), wobei \(\mathfrak A,\mathfrak B,\mathfrak C\) usw. verschiedene ungerade Primzahlen bezeichnen, so ist die Gesamtzahl aller eigentlich primitiven Formen in \(V,V',V''\) usw. \(=\frac{A\mathfrak{nabc}\ldots}{2\mathfrak{ABC}\ldots}\), wobei man setzen muss
\[
\begin{aligned}
\mathfrak n&=1\ (\text{wenn }\tfrac{4D}{AA}\equiv1,\ \mathrm{mod.}\,8),\text{ oder}\\
\mathfrak n&=2\ (\text{wenn }\tfrac{D}{AA}\text{ ganzzahlig ist}),\text{ oder}\\
\mathfrak n&=3\ (\text{wenn }\tfrac{4D}{AA}\equiv5,\ \mathrm{mod.}\,8);\text{ ferner}\\
\mathfrak a&=\mathfrak A\ (\text{wenn }\mathfrak A\text{ ein Teiler von }\tfrac{4D}{AA}\text{ ist}),\text{ oder}\\
\mathfrak a&=\mathfrak A\pm1\ (\text{wenn }\mathfrak A\text{ kein Teiler von }\tfrac{4D}{AA}\text{ ist,}\\ &\qquad\text{mit dem oberen Vorzeichen}\\
&\qquad\text{oder dem unteren, je nachdem, ob }\tfrac{4D}{AA}\text{ quadratischer Nichtrest}\\ &\qquad\text{oder quadratischer Rest modulo }\mathfrak A\text{ ist})
\end{aligned}
\]
und schließlich \(\mathfrak b,\mathfrak c\) usw. aus \(\mathfrak B,\mathfrak C\) ebenso ableiten muss wie \(\mathfrak a\) aus \(\mathfrak A\). Den Beweis hier ausführlicher darzustellen. erlaubt die Kürze nicht.

V. Was nun die Anzahl der Klassen betrifft, die durch die eigentlich primitiven Formen in \(V,V',V''\) usw. gegeben werden, so sind die folgenden drei Fälle zu unterscheiden.

\textit{Erstens}: Ist \(D\) eine negative Zahl, so bildet jede eigentlich primitive Form in \(V,V'\) usw. eine eigene Klasse; die Anzahl der gesuchten Klassen wird also durch die Formel der vorigen Bemerkung ausgedrückt, mit zwei Ausnahmen, nämlich wenn \(\frac{4D}{AA}\) entweder \(=-4\) oder \(=-3\) ist, also wenn \(D\) entweder \(=-AA\) oder \(=-\frac34 AA\) ist. Zum Beweis dieses Satzes muss offenbar nur gezeigt werden, dass zwei verschiedene Formen aus \(V,V',V''\) usw. nicht eigentlich äquivalent sein können. Wir nehmen also an, \((hh,i,k)\), \((h'h',i',k')\) seien zwei verschiedene eigentlich primitive Formen aus \(V,V',V''\) usw., die derselben Klasse angehören, und die erste gehe durch die eigentliche Substitution \(\alpha,\beta,\gamma,\delta\) in die zweite über. Dann haben wir die Gleichungen
\[
\begin{gathered}\alpha\delta-\beta\gamma=1,\\ hh\alpha\alpha+2i\alpha\gamma+k\gamma\gamma=h'h',\\ hh\alpha\beta+i(\alpha\delta+\beta\gamma)+k\gamma\delta=i'\end{gathered}
\]
Daraus folgt leicht, \textit{erstens}, dass \(\gamma\) gewiss nicht \(=0\) ist (woraus \(\alpha=\pm1\), \(hh=h'h'\), \(i'\equiv i\ (\mathrm{mod.}\,hh)\) folgen würden und damit die Identität der vorgelegten Formen, entgegen der
\SourceLeaf{287}
Annahme); \textit{zweitens}, dass \(\gamma\) durch den größten gemeinsamen Teiler der Zahlen \(h,h'\) teilbar ist; (setzt man nämlich diesen Teiler \(=r\), so teilt er offenbar auch \(2i,2i'\), ist aber zu \(k\) teilerfremd; ferner teilt \(rr\) die Zahl \(hhk-h'h'k'=ii-i'i'\), woraus leicht folgt, dass \(r\) auch \(i-i'\) teilt. Nun gilt aber \(\alpha i'-\beta h'h'=\alpha i+\gamma k\), also ist \(\gamma k\) und somit auch \(\gamma\) durch \(r\) teilbar); \textit{drittens}, dass \((\alpha hh+\gamma i)^2-D\gamma\gamma=hhh'h'\). Setzt man also \(\alpha hh+\gamma i=rp\), \(\gamma=rq\), so sind \(p\) und \(q\) ganze Zahlen, von denen die zweite nicht \(=0\) ist, und es gilt \(pp-Dqq=\frac{hhh'h'}{rr}\). Nun ist \(\frac{hhh'h'}{rr}\) die kleinste durch \(hh\) und \(h'h'\) zugleich teilbare Zahl und teilt somit \(AA\) und daher auch \(4D\). Also ist \(\frac{4Drr}{hhh'h'}\) eine (negative) ganze Zahl; setzt man sie \(=-e\), so gilt \(pp-Dqq=-\frac{4D}{e}\) oder \(4=(\frac{2rp}{hh'})^2+eqq\). In dieser Gleichung muss der Term \((\frac{2rp}{hh'})^2\), als Quadrat kleiner als \(4\), entweder \(0\) oder \(1\) sein. Im ersten Fall gilt \(eqq=4\) und \(D=-(\frac{hh'}{rq})^2\); daraus folgt, dass \(\frac{4D}{AA}\) ein mit negativem Vorzeichen versehenes Quadrat ist und somit gewiss nicht \(\equiv1\ (\mathrm{mod.}\,4)\), also \(O\) weder eine uneigentlich primitive noch eine von einer uneigentlich primitiven Ordnung abgeleitete Ordnung ist. Daher ist \(\frac{D}{AA}\) ganzzahlig, woraus leicht folgt, dass \(e\) durch \(4\) teilbar ist, dass \(qq=1\), \(D=-(\frac{hh'}{r})^2\) gilt und auch \(\frac{AA}{D}\) ganzzahlig ist. Somit muss \(D=-AA\) oder \(\frac{D}{AA}=-1\) gelten; dies ist die erste Ausnahme. Im zweiten Fall ist \(eqq=3\), also \(e=3\) und \(4D=-3(\frac{hh'}{r})^2\); daher ist \(3(\frac{hh'}{rA})^2\) eine ganze Zahl, die, da sie mit dem ganzzahligen Quadrat \((\frac{rA}{hh'})^2\) multipliziert \(3\) ergibt, nur \(3\) sein kann. Somit gilt \(4D=-3AA\) oder \(D=-\frac34 AA\); dies ist die zweite Ausnahme. In allen übrigen Fällen gehören also sämtliche eigentlich primitiven Formen in \(V,V',V''\) usw. zu verschiedenen Klassen. --- Für die Ausnahmefälle genügt es, die Ergebnisse einer nicht schwierigen Untersuchung anzugeben, die wir hier der Kürze wegen weglassen müssen. Im ersten Fall gehören von den eigentlich primitiven Formen in \(V,V',V''\) usw. stets je zwei, im zweiten je drei zu derselben Klasse; somit ist die Gesamtzahl aller gesuchten Klassen im ersten Fall die Hälfte, im zweiten ein Drittel des Wertes des in der vorigen Bemerkung angegebenen Ausdrucks.

\textit{Zweitens}: Ist \(D\) eine positive Quadratzahl, so bildet jede eigentlich primitive Form in \(V,V',V''\) usw. ausnahmslos eine eigene Klasse. Nehmen wir nämlich an, \((hh,i,k)\), \((h'h',i',k')\) seien zwei solche verschiedenen eigentlich äquivalenten Formen, und die erste gehe durch die eigentliche Substitution \(\alpha,\beta,\gamma,\delta\) in die zweite über. Dann gelten offenbar auch hier alle Überlegungen des vorigen Falls, in denen nicht vorausgesetzt wird, dass \(D\) negativ ist. Bezeichnen also \(p,q,r\)
\SourceLeaf{288}
dasselbe wie dort, so ist auch hier \(\frac{4Drr}{hhh'h'}\) ganzzahlig, nun aber positiv statt negativ und außerdem ein Quadrat. Setzt man es \(=gg\), so ist \((\frac{2rp}{hh'})^2-ggqq=4\), \textit{was widersinnig ist,} da die Differenz zweier Quadrate nicht \(4\) sein kann, wenn das kleinere Quadrat nicht \(0\) ist. Die Annahme kann daher nicht bestehen.

Für den \textit{dritten} Fall dagegen, in dem \(D\) eine positive Nichtquadratzahl ist, haben wir bisher keine allgemeine Regel, um die Anzahl der eigentlich primitiven Formen in \(V,V',V''\) usw. mit der Anzahl der daraus entstehenden verschiedenen Klassen zu vergleichen. Wir können allerdings sagen, dass die zweite Anzahl der ersten entweder gleich oder ein Teiler von ihr ist. Wir haben sogar einen besonderen Zusammenhang zwischen dem Quotienten dieser Zahlen und den kleinsten Werten von \(t,u\), die der Gleichung \(tt-Duu=AA\) genügen, entdeckt; ihn hier darzustellen würde aber zu weit führen. Ob sich jener Quotient in allen Fällen allein durch Betrachtung der Zahlen \(D,A\) erkennen lässt (wie in den vorigen Fällen), darüber können wir nichts Sicheres sagen. Hier sind einige Beispiele, deren Anzahl jeder leicht vermehren kann. Für \(D=13\), \(A=2\) ist die Anzahl der eigentlich primitiven Formen in \(V\) usw. gleich \(3\); alle sind äquivalent, bilden also nur eine Klasse. Für \(D=37\), \(A=2\) gibt es ebenfalls drei eigentlich primitive Formen in \(V\) usw., die zu drei verschiedenen Klassen gehören; für \(D=588\), \(A=7\) gibt es acht eigentlich primitive Formen in \(V\) usw., die vier Klassen bilden; für \(D=867\), \(A=17\) gibt es in \(V\) usw. \(18\) eigentlich primitive Formen, ebenso viele für \(D=1445\), \(A=17\), die aber bei der ersten Determinante in zwei, bei der zweiten in sechs Klassen zerfallen.

VI. Wendet man diese allgemeine Theorie auf den Fall an. in dem \(O\) eine uneigentlich primitive Ordnung ist, so folgt: Die Anzahl der Klassen in dieser Ordnung verhält sich zur Gesamtzahl aller Klassen in der eigentlich primitiven Ordnung wie \(1\) zur Anzahl der verschiedenen eigentlich primitiven Klassen, die durch die drei Formen \((1,0,-D)\), \((4,1,\frac{1-D}{4})\), \((4,3,\frac{9-D}{4})\) gebildet werden. Dabei entsteht eine einzige Klasse, wenn \(D\equiv1\ (\mathrm{mod.}\,8)\), weil in diesem Fall die zweite und dritte Form uneigentlich primitiv sind; gilt dagegen \(D\equiv5\ (\mathrm{mod.}\,8)\), so sind alle drei Formen eigentlich primitiv und ergeben ebenso viele verschiedene Klassen, wenn \(D\) negativ ist, mit der einzigen Ausnahme \(D=-3\), bei der sie eine einzige Klasse bilden. Schließlich gehört der Fall, dass \(D\) positiv ist (von der Form \(8n+5\)), zu denen, für die bisher eine allgemeine Regel fehlt. Wir können aber sagen, dass jene drei Formen in diesem Fall
\SourceLeaf{289}
entweder zu drei verschiedenen Klassen oder zu einer einzigen gehören, niemals zu zweien. Denn gehören die Formen \((1,0,-D)\), \((4,1,\frac{1-D}{4})\), \((4,3,\frac{9-D}{4})\) jeweils zu den Klassen \(K,K',K''\), so gilt offenbar \(K+K'=K'\), \(K'+K'=K''\); nimmt man also an, dass \(K\) und \(K'\) identisch sind, so sind auch \(K'\) und \(K''\) identisch. Nimmt man ebenso \(K\) und \(K''\) als identisch an, so sind es auch \(K'\) und \(K''\); da schließlich \(K'+K''=K\) gilt, folgt aus der Annahme, \(K'\) und \(K''\) seien identisch, dass auch \(K\) und \(K''\) zusammenfallen. Somit sind entweder alle drei Klassen \(K,K',K''\) verschieden oder alle drei identisch. \textit{Zum Beispiel} gibt es unterhalb von \(600\) \(75\) Zahlen der Form \(8n+5\), darunter \(16\) Determinanten, für die der erste Fall gilt, bei denen also die Anzahl der Klassen in der eigentlich primitiven Ordnung dreimal so groß ist wie in der uneigentlich primitiven, nämlich \(37,\allowbreak 101,\allowbreak 141,\allowbreak 189,\allowbreak 197,\allowbreak 269.\ 325,\allowbreak 333,\allowbreak 349,\allowbreak 373,\allowbreak 381,\allowbreak 389,\allowbreak 405,\allowbreak 485,\allowbreak 557,\allowbreak 573\); für die übrigen \(59\) gilt der zweite Fall, die Klassenzahl ist also in beiden Ordnungen gleich.

VII. Es ist kaum nötig zu bemerken, dass man durch die vorige Untersuchung nicht nur die Klassenzahlen verschiedener Ordnungen derselben Determinante vergleichen kann, sondern dass sie auch auf beliebige verschiedene Determinanten anwendbar ist, deren Verhältnis ein Verhältnis von Quadraten ist. Bezeichnet nämlich \(O\) irgendeine Ordnung der Determinante \(dmm\) und \(O'\) eine Ordnung der Determinante \(dm'm'\), so lässt sich \(O\) mit der eigentlich primitiven Ordnung der Determinante \(dmm\) vergleichen und diese mit der von der eigentlich primitiven Ordnung der Determinante \(d\) abgeleiteten Ordnung, oder, was hinsichtlich der Klassenzahl auf dasselbe hinausläuft, mit dieser Ordnung selbst; mit derselben lässt sich auf genau ähnliche Weise die Ordnung \(O'\) vergleichen.
